\chapter{Tests, validation et résultats}

\section{Stratégie}
La validation suit l'identité d'un événement à quatre niveaux : transaction PostgreSQL, tables Bronze, modèles Silver et tables Gold. Elle est complétée par le rafraîchissement Power BI. Le dépôt contient 107 tests dbt portant sur la non-nullité, l'unicité, les relations, les timestamps et les calculs métier.

\begin{figure}[htbp]
\centering
\resizebox{.9\textwidth}{!}{%
\begin{tikzpicture}[stage/.style={rounded corners,draw=walmartnavy,fill=softblue,minimum width=2.4cm,minimum height=.9cm,align=center},check/.style={diamond,aspect=2,draw=codegreen,fill=green!8,align=center,font=\scriptsize},arr/.style={-{Latex[length=2.3mm]},thick,walmartblue}]
\node[stage] (s) {Source}; \node[check,right=.8cm of s] (c1) {PK/FK}; \node[stage,right=.8cm of c1] (b) {Bronze}; \node[check,right=.8cm of b] (c2) {curseur};
\node[stage,below=1.2cm of c2] (si) {Silver}; \node[check,left=.8cm of si] (c3) {107 tests}; \node[stage,left=.8cm of c3] (g) {Gold}; \node[stage,left=.8cm of g] (p) {Power BI};
\draw[arr] (s)--(c1); \draw[arr] (c1)--(b); \draw[arr] (b)--(c2); \draw[arr] (c2)--(si); \draw[arr] (si)--(c3); \draw[arr] (c3)--(g); \draw[arr] (g)--(p);
\end{tikzpicture}}
\caption{Chaîne de validation d'un événement}
\label{fig:validation-chain}
\end{figure}

\section{Contrôles source}
Les requêtes suivantes doivent retourner zéro :

\begin{lstlisting}[language=SQL,caption={Contrôles d'intégrité PostgreSQL}]
select count(*)
from raw.order_items i
left join raw.orders o on o.order_id = i.order_id
where o.order_id is null;

select count(*)
from raw.orders o
left join raw.employees e on e.employee_id = o.employee_id
where e.employee_id is null or e.store_id <> o.store_id;
\end{lstlisting}

Un contrôle complémentaire compare le total d'en-tête à la somme des lignes. Un test d'erreur forcée doit montrer qu'aucun en-tête orphelin ne subsiste après rollback.

\section{Matrice de tests dbt}
\begin{table}[htbp]
\centering
\caption{Principales familles de tests}
\label{tab:test-matrix}
\begin{tabularx}{\textwidth}{>{\bfseries}p{3.4cm}p{4cm}X}
\toprule
Niveau & Exemple & Risque couvert \\
\midrule
Silver technique & Clé unique & Doublon après rejeu \\
Silver métier & Montant de ligne & Incohérence financière \\
Gold dimension & Une version courante & Doublon côté 1 de Power BI \\
Gold fait & Item unique et non nul & Grain instable \\
Intertables & Employé du magasin & Relation métier invalide \\
Régression & Item minimal par commande & Nullité des nouvelles commandes \\
\bottomrule
\end{tabularx}
\end{table}

Pour les commandes supérieures à 10 000, le test compare la référence de ligne de \tech{dim\_orders} au minimum calculé dans \tech{fact\_orders}. Il valide la règle exacte, et pas seulement la non-nullité.

\section{Test de bout en bout observé}
Un scénario exécuté durant le développement a suivi le client 2002, la commande 10015 et la ligne 30070. Les enregistrements ont été vus dans PostgreSQL, puis les onze tâches Airflow ont réussi. Les clés ont été retrouvées en Bronze, Silver et Gold, avant leur chargement dans Power BI. Ces identifiants documentent un cas de recette ; la volumétrie courante peut avoir augmenté depuis.

\begin{lstlisting}[language=SQL,caption={Recette Gold des nouvelles commandes}]
select order_id, order_item_id
from walmart.gold.dim_orders
where order_id > 10000
  and dbt_valid_to = timestamp('9999-12-31');

select f.order_id, min(f.order_item_id), d.order_item_id
from walmart.gold.fact_orders f
join walmart.gold.dim_orders d on d.order_id = f.order_id
where f.order_id > 10000
group by f.order_id, d.order_item_id
having min(f.order_item_id) <> d.order_item_id
    or d.order_item_id is null;
\end{lstlisting}

La seconde requête doit retourner zéro ligne. Un rafraîchissement Power BI est ensuite nécessaire : un import local ne se met pas à jour uniquement parce que Gold a changé.

\section{Résultats et incident analysé}
Le pipeline relie désormais les deux événements générés aux dimensions et au fait. L'historique des employés est complet, la relation 1:N peut être créée, et les commandes générées possèdent un \tech{order\_item\_id} conforme. La table de sauvegarde pré-correction est exclue du modèle BI.

Un incident Airflow a interrompu \tech{dbt run --select silver\_b} par \tech{SIGTERM} après que l'état de la task instance a été marqué en échec. Aucun message SQL dbt n'était présent avant l'arrêt. Les exécutions ultérieures ayant réussi, l'incident est interprété comme une interruption d'orchestration ou de synchronisation d'état. Le warning de sémaphores est une conséquence probable de l'arrêt brutal, pas une preuve de sa cause.

\section{Limites}
La latence analytique peut approcher une heure. Les données synthétiques ne reproduisent pas toutes les saisons et promotions. Le modèle ne possède pas encore une dimension date Gold dédiée. Les métriques de coût, délai P95 et disponibilité ne sont pas mesurées sur une période assez longue pour annoncer un SLA.

L'environnement Docker local n'est pas hautement disponible et les secrets ne sont pas centralisés. Une page de supervision devrait consolider la dernière version traitée, les volumes par couche, la durée du DAG et les erreurs de tests.

\keypoint{La réussite ne se limite pas à voir un nombre changer dans Power BI : elle consiste à prouver qu'un événement conserve ses clés, son grain et ses règles métier à travers chaque contrat de données.}

